% ===================================================================
%  1. TAULA — Eskola. — Lizeoa. — Ikasgela.
%
%  Ceci n'est PAS une transcription : c'est une traduction, faite en
%  2026, en basque unifie (euskara batua). D'ou trois differences avec
%  texto/io et texto/fr, et trois seulement :
%
%   * pas de \nl ni de \cc ;
%   * pas de \begin{VUpage} ;
%   * le gras se pose sur le TERME seul, ponctuation dehors.
%
%  Les cles « %%K » sont celles de texto/io, macro pour macro pour
%  l'apparat, et les renvois \textsuperscript{(n)} visent les MEMES
%  objets de la planche.
%
%  NEUVIEME COLONNE DES DIX-SEPT, ET LA PREMIERE DE CE GROUPE QUI NE
%  SOIT PAS INDO-EUROPEENNE. Le basque n'a de parent nulle part : ni le
%  finnois, qui est finno-ougrien, ni aucune des quinze autres ne lui
%  tend un mot. Toute colonne voisine cesse donc d'etre un temoin
%  utile, et il n'en reste qu'un — l'ido lui-meme, et le francais pour
%  verifier QUEL objet de la planche est vise.
%
%  IL DECLINE, ET LA DECLINAISON REMPLACE LA PREPOSITION, comme en
%  finnois : « ikasgelan » est « dans la classe », « mahaiaren gainean »
%  est « sur la table ». Le gras porte donc la forme FLECHIE, celle que
%  la phrase demande.
%
%  ET SON GENITIF SE PREPOSE : « gazteluaren atea », le chateau
%  d'abord, la porte ensuite. On attend donc des inversions, comme en
%  finnois et pour la meme raison. La ou le finnois s'en tirait par une
%  relative, le basque ne le peut pas toujours : sa relative se prepose
%  elle aussi, et deplacerait le renvoi du meme cote. Le remede sera
%  donc l'apposition postposee — « atea (8), gaztelu zaharrarena (7) »
%  — qui est du basque courant et rend l'ordre de l'ido.
%
%  ORDRE DES MOTS : le verbe se met a la fin, et ce qui le suit est
%  arriere-plan. C'est justement ce qui sauve le tableau 1 :
%  « Irakaslea (7) aulki batean (8) eserita dago katedraren (6)
%  atzean » donne 7, 8, 6 comme l'ido, la ou l'ordre neutre aurait
%  donne 7, 6, 8.
%
%  « postigo » (t01, renvoi 78) EST UN VASISTAS, ET LA PLANCHE LE DIT.
%  Le mot a ete verifie sur la gravure en 2026, non sur les
%  dictionnaires : la verriere de la classe a deux rangees de carreaux,
%  et dans la rangee HAUTE deux petits vantaux sont graves battants,
%  gonds sur un montant vertical, ouverts vers l'interieur. C'est le
%  seul ouvrant de la verriere, et il ne sert qu'a aerer.
%
%  D'OU CE QU'IL N'EST PAS, et ou seize colonnes s'etaient trompees :
%  ni une lucarne de toit (it. « lucernari »), ni un volet (cs, lt, lb),
%  ni un meneau ni un montant de chassis (ar, ca, es, oc), ni une
%  imposte fixe, ni le battant de la grande fenetre (bn), ni le
%  « ranma » de la cloison japonaise. Et « framuga » russe nomme
%  l'imposte entiere, quand le francais nomme le VANTAIL : c'est
%  « fortochka » qu'il fallait, comme l'ukrainien avait deja
%  « kvatyrka ».
% ===================================================================

%%K t01-apar-0 apar
\VUsaut{2.0mm}
\VUcentre{12.6pt}{120}{AZALPEN-LIBURUXKA}
\VUsaut{3.2mm}
\VUcentre{8.4pt}{160}{DELMASEN}
\VUsaut{2.6mm}
\VUtitre{15.6pt}{0}{laguntza-taulentzat}
\VUsaut{3.4mm}
\VUornamento{9pt}
\VUsaut{5.0mm}
\VUcentre{13.2pt}{90}{LEHEN SAILA}
\VUsaut{3.0mm}
\VUfilet{16mm}
\VUsaut{5.6mm}

%%K t01-apar-1 apar
\VUtitre{15.0pt}{60}{1. TAULA}
\VUsaut{1.4mm}
\VUtitre{11.4pt}{0}{Eskola. --- Lizeoa. --- Ikasgela.}
\VUsaut{2.2mm}
\VUfilet{20mm}
\VUsaut{6.4mm}

%%K t01-tit-1 sub
\VUpk{10.2pt}{40}{Deskribapen orokorra.}
\VUsaut{3.0mm}

%%K t01-01-1 p
1. --- Taula honetan lizeoko \VUgras{ikasgela} bat ageri da. Ikasgela honetan
zortzi \VUgras{mahai} \textsuperscript{(18)} eta zortzi \VUgras{banku}
\textsuperscript{(32)} daude. Banku bakoitzean hiru ikasle eserita daude.
\VUgras{Irakaslea} \textsuperscript{(7)} \VUgras{aulki} batean
\textsuperscript{(8)} eserita dago bere \VUgras{katedraren} \textsuperscript{(6)}
atzean.

%%K t01-02-1 p
2. --- Katedraren ezkerrean beirazko \VUgras{armairu} bat \textsuperscript{(15)}
dago. Eskuinean, \VUgras{arbela} \textsuperscript{(1)}, bere \VUgras{trapuarekin}
\textsuperscript{(2)} eta \VUgras{klarionarekin} \textsuperscript{(3)}.

%%K t01-02-2 p
Katedraren gainean \VUgras{horma-irudiak} \textsuperscript{(9, 11, 12)} eta
\VUgras{mapa} bat \textsuperscript{(10)} zintzilik daude.

%%K t01-03-1 p
3. --- Ikasle guztiak ez daude beren \VUgras{lekuetan} \textsuperscript{(17)}.
Joanes \textsuperscript{(4)} arbelaren ondoan zutik dago; trapua eskuan darama eta
\VUgras{geometria-irudiak} ezabatzen ari da.

%%K t01-03-2 p
Petri katedraren ondoan zutik dago; \VUgras{idazlanak} \textsuperscript{(5)}
biltzen ari da irakasleari eramateko.

%%K t01-04-1 p
4. --- \VUgras{Irakasle} hau \VUgras{hizkuntza modernoetakoa} da. Ikasleei
atzerriko hizkuntzetan hitz egiten, irakurtzen eta idazten erakusten die
\VUgras{(alemana, ingelesa, gaztelania, italiera, errusiera} edo
\VUgras{frantsesa)}.

%%K t01-05-1 p
5. --- Ikasgela oso handia eta oso argitsua da. Barrenean \VUgras{leiho} zabal
bat dago, bi \VUgras{aireztapen-leihorekin} \textsuperscript{(78)}, eta \VUgras{beirazko
ate} bat, gela aireztatzeko eta argitzeko.

%%K t01-05-2 p
Ikasgela honetan ez da hotzik egiten. Erdian, gela berotzeko, \VUgras{berogailu}
oso on bat \textsuperscript{(46)} dago, bere \VUgras{ke-hodiarekin}
\textsuperscript{(45)}.

%%K t01-05-3 p
Trenkadan \VUgras{esekitokiak} \textsuperscript{(58)} iltzatuta daude; haietan
ikasleen kapelak eta berokiak zintzilikatzen dira.

%%K t01-tit-2 sub
\VUsaut{5.2mm}
\VUpk{10.2pt}{40}{Jolastokia.}
\VUsaut{3.6mm}

%%K t01-06-1 p
6. --- \VUgras{Erlojuak} \textsuperscript{(63)} bederatziak eta bost markatzen
ditu.

%%K t01-06-2 p
\VUgras{Jolasaldia} \textsuperscript{(76)} amaitu berri da, eta ikasgaia berriro
hasten ari da. Jolasaldia \VUgras{jolastokian} \textsuperscript{(75)} igarotzen
da. Zenbait ikasle jolasean ikusten ditugu. \VUgras{Zaintzaileak}
\textsuperscript{(74)} begiratzen die.

%%K t01-07-1 p
7. --- Jolastokian beste jende bat ere ikusten dugu: \VUgras{ikuskatzailea}
\textsuperscript{(73)}, \VUgras{zuzendaria} \textsuperscript{(71)} eta
\VUgras{ikasketaburua} \textsuperscript{(72)}.

%%K t01-07-2 p
Ikuskatzaileak eskolak aztertzen ditu, zuzendariak ikastetxea gobernatzen du,
ikasketaburuak ikasleen lanari begiratzen dio.

%%K t01-07-3 p
Laugarrena \VUgras{atezaina} \textsuperscript{(77)} da. Erregistroa besapean
darama, faltak markatzeko.

%%K t01-08-1 p
8. --- Jolastokiaren barrenean \VUgras{gimnasia-tresnak} \textsuperscript{(70)}
eta \VUgras{eskulanaren} \textsuperscript{(69)} lantegia daude.

%%K t01-08-2 p
Taularen eskuinean ozta-ozta ikusten dira \VUgras{musikako} eta
\VUgras{marrazketako} \VUgras{ikasgelak}.

%%K t01-noto-1 noto
\VUnotes{143.2mm}{%
(*) \textit{Izenentzat} ere latinezko forman ematea gomendatzen dugu. Horrela
bakarrik saihets daitezke kontaezinezko aldaerak. Eta nola hautatu
\textit{Giovanni, Jean, Johann, Jan, Hans, John, Joanes} eta gainerakoen
artean? Ez dugu izenen hiztegi berezirik egingo. Har dezagun latinezko
nominatiboa eta esan dezagun: \VUgras{Ioannes, Ioanna; Iulius, Iulia; Iakobus;
Andreas; Lukas.} (Gramatika osoa).%
}

%%K t01-tit-3 sub
\VUpk{10.2pt}{40}{Deskribapen zehatza.}
\VUsaut{2.4mm}

%%K t01-tit-4 sub
\VUpk{10.2pt}{40}{Ikasle onak.}
\VUsaut{3.0mm}

%%K t01-09-1 p
9. --- Katedraren eskuinean, lehen bankuko ikasleak \VUgras{Henrike} (*),
\VUgras{Estepan} eta \VUgras{Karlos} dira.

%%K t01-09-2 p
Henrike oso adi eta oso saiatua da; irakasleari entzuten dio. Mahai gainean
\VUgras{koaderno} bat \textsuperscript{(28)}, \VUgras{liburu} bat
\textsuperscript{(29)}, \VUgras{hiztegi} bat \textsuperscript{(30)} eta
\VUgras{arkatz-ontzi} bat \textsuperscript{(31)} ditu. Bere liburuak
\VUgras{orrialde} asko ditu; kapitulu bakoitzean \VUgras{paragrafo} batzuk daude,
eta \VUgras{lerro} bakoitzean \VUgras{hitz} asko.

%%K t01-09-4 p
Estepan bere aldamenekoa \textsuperscript{(24)} langilea da. Hurrengorako emandako
ikasgaia koadernoan idazten ari da. \VUgras{Idazkera} bikaina du. Bere liburua
itxita dago. \VUgras{Azala} \textsuperscript{(27)} baino ez zaio ikusten. Haren
ondoan bere \VUgras{konpasa} \textsuperscript{(26)} datza.

%%K t01-09-5 p
Karlosek \textsuperscript{(23)} liburua eskuan dauka; irekitzen ari da, ikasgaia
berriro errepasatzeko. Ikasle orori bezala, aurrean \VUgras{tintontzi} bat
\textsuperscript{(25)} du. Esanekoa da.

%%K t01-10-1 p
10. --- Hurrengo mahaian \VUgras{Ludobiko, Anton} eta \VUgras{Paulo} daude
eserita. Ludobikok bere \VUgras{ikasgaia} \textsuperscript{(22)} errezitatzen du
eta irakaslearen \VUgras{galderei} erantzuten die. Anton \textsuperscript{(21)}
oso adigabea da; batzuetan irakasleak zigortu ere egiten du. Paulo
\textsuperscript{(20)} lagun ona da, atsegina eta prestua. Oso adi dago.

%%K t01-tit-5 sub
\VUsaut{4.2mm}
\VUpk{10.2pt}{40}{Ikasle txarrak.}
\VUsaut{3.2mm}

%%K t01-11-1 p
11. --- Henrikeren atzean \VUgras{Jurgi} \textsuperscript{(38)} dago eserita. Ez
da mutil gaiztoa, baina \VUgras{berritsua} da, adigabea eta ezin geldi egon
daitekeena. Bere \VUgras{arinkeriak} eta \VUgras{desobedientziak}
\VUgras{nahastea} sartzen dute ikasgaian, eta sarritan \VUgras{ohartarazpen}
zorrotza merezi du. Bere \VUgras{zirriborroa} \textsuperscript{(35)} zikina dago;
\VUgras{tinta-orbanak} egiten ditu bere \VUgras{lumaz} \textsuperscript{(34)},
eta horregatik etengabe darabil bere \VUgras{goma} \textsuperscript{(36)}; bere
\VUgras{motxila} \textsuperscript{(33)} urratuta dago, hain da zabarra. Zertarako
dakar bere \VUgras{idazluma} \textsuperscript{(37)} hizkuntza modernoen
ikasgelara?

%%K t01-11-2 p
Bere aldamenekoak, Edmundek \textsuperscript{(39)}, ez du maite. Egia da pixka bat
alferra dela, baina isila eta lasaia da. Ez da berritsua; entzun beharrean,
\VUgras{ipuinak} irakurtzea nahiago du bere \VUgras{irakurgai-liburuan}.

%%K t01-11-3 p
Banku honetako hirugarren ikaslea \VUgras{Ludwig} \textsuperscript{(40)} da.
Saiatua da. \VUgras{Paper-orri} zuri baten gainean idazten du. Aurrean
\VUgras{xurgapapera} \textsuperscript{(41)} jarri du.

%%K t01-12-1 p
12. --- Irakaslearen ezkerrean, bigarren bankuko ikasleak oso txikiak dira.
Lehena, \VUgras{Emil} txikia, ez da oraindik ondo idazten trebea; bere
\VUgras{arbeltxoa} \textsuperscript{(42)}, bere \VUgras{arbel-arkatza} eta bere
\VUgras{belakia} \textsuperscript{(43)} dakartza.

%%K t01-12-2 p
Haren ondoan \VUgras{Augusto} eta \VUgras{Bernardo} \textsuperscript{(44)} daude.
Azken honek ondo lan egiten du, baina \VUgras{itxura} oso zabarra du.

%%K t01-13-1 p
13. --- Haien atzean hiru \VUgras{alfer} daude eserita. \VUgras{Albertok} ez ditu
ikasgaiak ikasten. \VUgras{Jakobe} \textsuperscript{(47)} lo dago entzun
beharrean. Bere \VUgras{alferkeriak} \VUgras{errietak} eta \VUgras{zigorrak} ere
ekartzen dizkio. Azkenik, \VUgras{Kristian} \textsuperscript{(48)} arinegia da.
Gaizki \VUgras{portatzen} da eta ez da zuzentzen saiatzen.

%%K t01-14-1 p
14. --- Haien aldamenekoak, aitzitik, ondo portatzen dira. \VUgras{Ernestok}
\textsuperscript{(51)} ordena eta zehaztasuna maite ditu; beti darabiltza bere
\VUgras{erregela} \textsuperscript{(50)} eta bere \VUgras{arkatza}
\textsuperscript{(49)}. \VUgras{Kanpo-ikaslea} da.

%%K t01-14-2 p
\VUgras{Frantzisko} \textsuperscript{(52)} \VUgras{barne-ikaslea} da; uniformea
darama, ikastetxean jaten eta lo egiten du. \VUgras{Lagun} ona da.

%%K t01-14-3 p
Mauriziok bere \VUgras{arkatza} zorrozten du bere \VUgras{kanibetaz}
\textsuperscript{(56)}; oso zuhurra da.

%%K t01-14-4 p
Bere liburu guztiak dakartza, bere \VUgras{gramatika} \textsuperscript{(55)},
bere \VUgras{ohar-koadernoa} \textsuperscript{(54)} eta bere \VUgras{ohar-bloka}
ere \textsuperscript{(53)}. Bere \VUgras{kartera} \textsuperscript{(57)} lurrean
datza, haren atzean.

%%K t01-14-5 p
Ikasgelaren barreneko bi mahaiak \VUgras{ikasle onek} \textsuperscript{(19)}
betetzen dituzte.

%%K t01-tit-6 sub
\VUsaut{4.6mm}
\VUpk{10.2pt}{40}{Marrazketa eta musika.}
\VUsaut{3.4mm}

%%K t01-15-1 p
15. --- \VUgras{Marrazketa-gelan} horman \VUgras{eredu} ederrak daude zintzilik,
eta sabaitik \VUgras{lanpara} bat \textsuperscript{(62)} dilindan dago, bere
\VUgras{pantailarekin}. \VUgras{Irakaslea} \textsuperscript{(60)} oso
pazientziatsua da; ikasleen \VUgras{marrazkiak} \textsuperscript{(61)} zuzentzen
ditu eta \VUgras{bustoa} \textsuperscript{(59)} naturaletik marrazten erakusten
die.

%%K t01-16-1 p
16. --- \VUgras{Musika-gela} oso atsegina dirudi. Han \VUgras{notak} irakurtzen,
\VUgras{eskala-mailak} \textsuperscript{(67)} kantatzen --- \VUgras{pentagraman}
idatziak (do, re, mi, fa, sol, la, si\,=\,C. D. E. F. G. A. B.) ---,
\VUgras{giltzak} ezagutzen eta \VUgras{doinu} zoragarriak kantatzen ikasten da.
Notak \VUgras{atrilaren} gainean \textsuperscript{(65)} jartzen dira. Ikasleetako
batek \VUgras{pianoa jotzen du} \textsuperscript{(64)}, eta
\VUgras{musika-irakasleak} \textsuperscript{(66)} \VUgras{konpasa markatzen du}
\textsuperscript{(68)}.

%%K t01-17-1 p
17. --- Ozta-ozta ikusten da \VUgras{eskulanaren} \textsuperscript{(69)}
lantegiaren izkina, non haurrek zura lantzen eta burdina limatzen ikas
dezaketen. Ez dago eskola guztietan.

%%K t01-tit-7 sub
\VUsaut{5.0mm}
\VUpk{10.2pt}{40}{Natur zientzien ikasgela.}
\VUsaut{3.6mm}

%%K t01-18-1 p
18. --- Matematikako irakasleak \VUgras{geometria,} \VUgras{aritmetika, kontuak}
eta \VUgras{sistema metrikoa} irakasten dizkie ikasleei. Taulan geometriako
irudi nagusiak ikusten ditugu: \VUgras{zirkulua} \textsuperscript{(a)},
\VUgras{karratua} \textsuperscript{(b)}, \VUgras{triangelua}
\textsuperscript{(c)}, \VUgras{marra} \textsuperscript{(d)}.

%%K t01-18-2 p
\VUgras{Eragiketa} bat ere ikusten dugu; ez da \VUgras{batuketa}, ez
\VUgras{kenketa}, ez \VUgras{biderketa}: \VUgras{zatiketa}
\textsuperscript{(e)} da.

%%K t01-19-1 p
19. --- Sistema metrikoaren taulan hauek ikusten ditugu: \VUgras{metroa}
\textsuperscript{(a)}, \VUgras{balantza} \textsuperscript{(b)} eta bere
\VUgras{pisuak}, \VUgras{litroa} \textsuperscript{(e)}, \VUgras{dekalitroa}
\textsuperscript{(f)}, \VUgras{dezilitroa} eta \VUgras{metro kubikoa}
\textsuperscript{(d)}.

%%K t01-19-2 p
Ikasleek \VUgras{natur zientziak} \textsuperscript{(11)} ere ikasten dituzte ---
zoologia \textsuperscript{(a)}, botanika \textsuperscript{(b)}, geologia
\textsuperscript{(c)} --- eta \VUgras{historia} \textsuperscript{(12)}.

%%K t01-19-3 p
Armairuan \VUgras{fisikarako} \textsuperscript{(15)} eta \VUgras{kimikarako}
\textsuperscript{(16)} tresnak daude.

%%K t01-19-4 p
Armairuaren gainean hauek daude: \VUgras{esfera} \textsuperscript{(14)},
\VUgras{konoa} eta \VUgras{globoa} \textsuperscript{(13)}.

%%K t01-19-5 p
Taulen gainean \VUgras{mapa} bat dago zintzilik, munduaren bost zatiak
erakusten dituena: \VUgras{Europa} \textsuperscript{(g)}, \VUgras{Asia}
\textsuperscript{(e)}, \VUgras{Afrika} \textsuperscript{(f)}, \VUgras{Amerika}
\textsuperscript{(ab)} eta \VUgras{Ozeania} \textsuperscript{(h)}, eta bi ozeano
handiak --- \VUgras{Barea} \textsuperscript{(c)} eta \VUgras{Atlantikoa}
\textsuperscript{(d)}.
\VUsaut{6.0mm}
\VUfilet{22mm}
